\documentclass{article}
\usepackage[T1]{fontenc}
\usepackage{lmodern}
\usepackage{graphicx}
\usepackage{amsmath,amssymb}
\usepackage[a4paper,margin=2cm]{geometry}
\usepackage[absolute,overlay]{textpos}
\setlength{\TPHorizModule}{1mm}
\setlength{\TPVertModule}{1mm}
\usepackage[indonesian]{babel}
\usepackage{hyperref}
\setlength{\emergencystretch}{2em}
% unit: Halaman 30

\begin{document}

% === Halaman 30 (python/native) ===

Penjumlahan Dua Titik di dalam EC pada GF(p)

Misalkan P(xp, yp) dan Q(xq, yq).  

Penjumlahan: P + Q = R

Koordinat Titik R: 

     xr = m2 – xp – xq mod p

     yr = m(xp – xr) – yp  mod p

 m adalah gradien:


30

𝑚= 𝑦𝑝−𝑦𝑞

𝑥𝑝−𝑥𝑞

mod 𝑝

\end{document}
